% Fibonacci in Rust
% Exported from Grok by ThreadPress on 2026-10-24.
% Compile with pdfLaTeX, XeLaTeX or LuaLaTeX.
\documentclass[11pt,a4paper]{article}
\usepackage{iftex}
\ifPDFTeX
  \usepackage[utf8]{inputenc}
  \usepackage[T1]{fontenc}
\else
  \usepackage{fontspec}
\fi
\usepackage[margin=2.2cm]{geometry}
\usepackage{amsmath, amssymb}
\usepackage[normalem]{ulem}
\usepackage{xcolor}
\usepackage{listings}
\usepackage{array}
\usepackage{graphicx}
\usepackage{enumitem}
\usepackage{hyperref}
\lstset{basicstyle=\ttfamily\small, breaklines=true, columns=fullflexible, keepspaces=true,
  frame=single, rulecolor=\color{black!15}, backgroundcolor=\color{black!3}, upquote=true}
\setlength{\parindent}{0pt}
\setlength{\parskip}{0.6em}
\title{Fibonacci in Rust}
\date{Grok \textperiodcentered{} 2026-10-24}
\begin{document}
\maketitle

\section*{You}

Write fibonacci in Rust


\section*{Grok}

Here is an \emph{iterative} version:

\begin{lstlisting}
fn fib(n: u64) -> u64 {
    let (mut a, mut b) = (0, 1);
    for _ in 0..n { (a, b) = (b, a + b); }
    a
}
\end{lstlisting}

\begin{enumerate}[start=3]
  \item Third point
  \item Fourth point
\end{enumerate}


\end{document}
